% GENERATED FILE. Edit canonical lessons, then run npm run generate:books.
% canonical-source-hash: fnv1a64:6eac4a729080c036
% canonical-lessons: GU-C239-dedko, GU-C239-kadav, GU-C239-dubvu, GU-C239-dankhvu, GU-C239-chatakvu

\chapter{Down by the pond}
\label{ch:gu-down-by-the-pond}
\hlchaptermodality{239}

\begin{chapteropening}
\textbf{By the end of this chapter:} \emph{I can say what happened down by the pond.}

Complete the last lesson of chapter 239: I can say what happened down by the pond.
\end{chapteropening}

\section[\texorpdfstring{\gu{દેડકો} (deḍko)}{દેડકો (deḍko)}]{\gu{દેડકો} (deḍko) — frog}

\label{lesson:GU-C239-dedko}



\begin{quote}
Before the new one: say the Gujarati for ant, then the Gujarati for to hang.
\end{quote}

\subsection*{You'll want to know: \gu{દેડકો}}
\textbf{\gu{દેડકો}} — \emph{deḍko} — \textquotedblleft{}frog\textquotedblright{}.

\begin{quote}
\textbf{\gu{રાત્રે} \gu{તળાવમાં} \gu{દેડકો} \gu{બોલ્યો}.} — \emph{rātre taḷāvmā̃ deḍko bolyo} — \textquotedblleft{}A frog croaked in the pond at night.\textquotedblright{}
\end{quote}

\begin{cousinweb}
\textbf{\gu{દેડકો}} (\emph{deḍko}) is a tadbhava, a word inherited from Sanskrit and reshaped by sound change. It comes from \emph{dardura}, \textquotedblleft{}frog\textquotedblright{} (CDIAL 6198).
\end{cousinweb}

\begin{grammarlens}[title={What you've built}]
One more word for everyday life at A2.
\end{grammarlens}

\subsection*{Your turn}
\begin{itemize}
\raggedright
  \item \emph{Say it:} \emph{deḍko}
  \item \emph{Say it:} \emph{deḍko}, once more
  \item \emph{Say it:} \emph{rātre taḷāvmā̃ deḍko bolyo} — A frog croaked in the pond at night.
  \item \emph{Recall:} say the Gujarati for ant, then the Gujarati for to hang, then say \emph{deḍko} again
\end{itemize}

\begin{tcolorbox}[breakable,colback=teal!4,colframe=teal!35!black,title={Before you move on}]
What does \gu{દેડકો} mean? (\textquotedblleft{}Frog\textquotedblright{}.) Say it once more.
\end{tcolorbox}
\section[\texorpdfstring{\gu{કાદવ} (kādav)}{કાદવ (kādav)}]{\gu{કાદવ} (kādav) — mud, mire}

\label{lesson:GU-C239-kadav}



\begin{quote}
Before the new one: say the Gujarati for to hang, then the Gujarati for frog.
\end{quote}

\subsection*{You'll want to know: \gu{કાદવ}}
\textbf{\gu{કાદવ}} — \emph{kādav} — \textquotedblleft{}mud, mire\textquotedblright{}.

\begin{quote}
\textbf{\gu{તળાવ} \gu{પાસે} \gu{ઘણો} \gu{કાદવ} \gu{હતો}.} — \emph{taḷāv pāse ghaṇo kādav hato} — \textquotedblleft{}There was a lot of mud near the pond.\textquotedblright{}
\end{quote}

\begin{cousinweb}
\textbf{\gu{કાદવ}} (\emph{kādav}) is a tadbhava, a word inherited from Sanskrit and reshaped by sound change. It comes from \emph{kardama}, \textquotedblleft{}mud, mire\textquotedblright{} (CDIAL 2869).
\end{cousinweb}

\begin{grammarlens}[title={What you've built}]
One more word for everyday life at A2.
\end{grammarlens}

\subsection*{Your turn}
\begin{itemize}
\raggedright
  \item \emph{Say it:} \emph{kādav}
  \item \emph{Say it:} \emph{kādav}, once more
  \item \emph{Say it:} \emph{taḷāv pāse ghaṇo kādav hato} — There was a lot of mud near the pond.
  \item \emph{Recall:} say the Gujarati for to hang, then the Gujarati for frog, then say \emph{kādav} again
\end{itemize}

\begin{tcolorbox}[breakable,colback=teal!4,colframe=teal!35!black,title={Before you move on}]
What does \gu{કાદવ} mean? (\textquotedblleft{}Mud, mire\textquotedblright{}.) Say it once more.
\end{tcolorbox}
\section[\texorpdfstring{\gu{ડૂબવું} (ḍūbvũ)}{ડૂબવું (ḍūbvũ)}]{\gu{ડૂબવું} (ḍūbvũ) — to sink, to drown}

\label{lesson:GU-C239-dubvu}



\begin{quote}
Before the new one: say the Gujarati for frog, then the Gujarati for mud.
\end{quote}

\subsection*{You'll want to know: \gu{ડૂબવું}}
\textbf{\gu{ડૂબવું}} — \emph{ḍūbvũ} — \textquotedblleft{}to sink, to drown\textquotedblright{}.

\begin{quote}
\textbf{\gu{મારો} \gu{બૂટ} \gu{કાદવમાં} \gu{ડૂબ્યો}.} — \emph{māro būṭ kādavmā̃ ḍūbyo} — \textquotedblleft{}My shoe sank into the mud.\textquotedblright{}
\end{quote}

\begin{cousinweb}
\textbf{\gu{ડૂબવું}} (\emph{ḍūbvũ}) comes from \emph{*ḍubb}, \textquotedblleft{}sink\textquotedblright{}, a reconstructed form that CDIAL reads as a reshuffled Middle Indo-Aryan \emph{buḍḍaï}; it glosses the Gujarati verb \textquotedblleft{}to sink\textquotedblright{} (CDIAL 5561).
\end{cousinweb}

\begin{grammarlens}[title={What you've built}]
One more word for everyday life at A2.
\end{grammarlens}

\subsection*{Your turn}
\begin{itemize}
\raggedright
  \item \emph{Say it:} \emph{ḍūbvũ}
  \item \emph{Say it:} \emph{ḍūbvũ}, once more
  \item \emph{Say it:} \emph{māro būṭ kādavmā̃ ḍūbyo} — My shoe sank into the mud.
  \item \emph{Recall:} say the Gujarati for frog, then the Gujarati for mud, then say \emph{ḍūbvũ} again
\end{itemize}

\begin{tcolorbox}[breakable,colback=teal!4,colframe=teal!35!black,title={Before you move on}]
What does \gu{ડૂબવું} mean? (\textquotedblleft{}To sink, to drown\textquotedblright{}.) Say it once more.
\end{tcolorbox}
\section[\texorpdfstring{\gu{ડંખવું} (ḍankhvũ)}{ડંખવું (ḍankhvũ)}]{\gu{ડંખવું} (ḍankhvũ) — to sting}

\label{lesson:GU-C239-dankhvu}



\begin{quote}
Before the new one: say the Gujarati for mud, then the Gujarati for to sink.
\end{quote}

\subsection*{You'll want to know: \gu{ડંખવું}}
\textbf{\gu{ડંખવું}} — \emph{ḍankhvũ} — \textquotedblleft{}to sting\textquotedblright{}.

\begin{quote}
\textbf{\gu{મને} \gu{મધમાખી} \gu{ડંખી}.} — \emph{mane madhmākhī ḍankhī} — \textquotedblleft{}A bee stung me.\textquotedblright{}
\end{quote}

\begin{cousinweb}
\textbf{\gu{ડંખવું}} (\emph{ḍankhvũ}) is built on \emph{ḍankh}, \textquotedblleft{}a sting\textquotedblright{}. It is usually linked to Sanskrit \emph{daṁśa}, \textquotedblleft{}stinging insect\textquotedblright{}, where CDIAL lists Gujarati \emph{ḍā̃s}, \textquotedblleft{}gadfly\textquotedblright{} (CDIAL 6110), but CDIAL does not list \emph{ḍankh}, so the link is not settled.
\end{cousinweb}

\begin{grammarlens}[title={What you've built}]
One more word for everyday life at A2.
\end{grammarlens}

\subsection*{Your turn}
\begin{itemize}
\raggedright
  \item \emph{Say it:} \emph{ḍankhvũ}
  \item \emph{Say it:} \emph{ḍankhvũ}, once more
  \item \emph{Say it:} \emph{mane madhmākhī ḍankhī} — A bee stung me.
  \item \emph{Recall:} say the Gujarati for mud, then the Gujarati for to sink, then say \emph{ḍankhvũ} again
\end{itemize}

\begin{tcolorbox}[breakable,colback=teal!4,colframe=teal!35!black,title={Before you move on}]
What does \gu{ડંખવું} mean? (\textquotedblleft{}To sting\textquotedblright{}.) Say it once more.
\end{tcolorbox}
\section[\texorpdfstring{\gu{છટકવું} (chaṭakvũ)}{છટકવું (chaṭakvũ)}]{\gu{છટકવું} (chaṭakvũ) — to slip away, to escape}

\label{lesson:GU-C239-chatakvu}



\begin{quote}
Before the new one: say the Gujarati for to sink, then the Gujarati for to sting.
\end{quote}

\subsection*{You'll want to know: \gu{છટકવું}}
\textbf{\gu{છટકવું}} — \emph{chaṭakvũ} — \textquotedblleft{}to slip away, to escape\textquotedblright{}.

\begin{quote}
\textbf{\gu{દેડકો} \gu{મારા} \gu{હાથથી} \gu{છટક્યો}.} — \emph{deḍko mārā hāththī chaṭakyo} — \textquotedblleft{}The frog slipped out of my hand.\textquotedblright{}
\end{quote}

\begin{cousinweb}
\textbf{\gu{છટકવું}} (\emph{chaṭakvũ}) comes from \emph{*chaṭṭ}, \textquotedblleft{}scatter\textquotedblright{}, a reconstructed form; CDIAL glosses the Gujarati verb \textquotedblleft{}to slip off\textquotedblright{} (CDIAL 4968).
\end{cousinweb}

\begin{grammarlens}[title={What you've built}]
One more word for everyday life at A2.
\end{grammarlens}

\subsection*{Your turn}
\begin{itemize}
\raggedright
  \item \emph{Say it:} \emph{chaṭakvũ}
  \item \emph{Say it:} \emph{chaṭakvũ}, once more
  \item \emph{Say it:} \emph{deḍko mārā hāththī chaṭakyo} — The frog slipped out of my hand.
  \item \emph{Recall:} say the Gujarati for to sink, then the Gujarati for to sting, then say \emph{chaṭakvũ} again
\end{itemize}

\begin{tcolorbox}[breakable,colback=teal!4,colframe=teal!35!black,title={Before you move on}]
What does \gu{છટકવું} mean? (\textquotedblleft{}To slip away, to escape\textquotedblright{}.) Say it once more.
\end{tcolorbox}
